\documentclass[11pt]{article}
\pdfinfoomitdate=1
\pdftrailerid{}
\pdfsuppressptexinfo=15
\usepackage[a4paper,margin=0.94in]{geometry}
\usepackage[T1]{fontenc}
\usepackage{lmodern}
\usepackage{amsmath,amssymb,amsthm,booktabs,microtype}
\usepackage[numbers,sort&compress]{natbib}
\usepackage{cleveref}

\newtheorem{theorem}{Theorem}[section]
\newtheorem{corollary}[theorem]{Corollary}
\newtheorem{proposition}[theorem]{Proposition}
\theoremstyle{remark}
\newtheorem{remark}[theorem]{Remark}
\newcommand{\cH}{\mathcal H}
\newcommand{\cY}{\mathcal Y}
\newcommand{\one}{\mathbf 1}

\title{A Sharp Post-First-Alias\\Mesoscopic Crossover}
\author{Bin Wang}
\date{August 2026}

\begin{document}
\maketitle

\begin{abstract}
We resolve the crossover between the first counterloop alias and the first
$L$ even coordinates immediately above it.  Let
$y_k=(\beta_kR)^2$, $x=(\beta R)^2>1$, and define
\[
 A_k=\frac{|s_{k,2k}|R^{2k}}{2k},\qquad
 B_k(L)=\sum_{j=1}^L
 \frac{|s_{k,2k+2j}|R^{2k+2j}}{2k+2j}.
\]
The exact graded ledger gives
\[
 A_k=(1-k^{-1})y_k^k,\qquad
 \frac{B_k(L)}{A_k}
 =\frac{k}{k-1}\sum_{j=1}^L\frac{y_k^j}{k+j}.
\]
For every integer envelope $\ell_k=o(k)$, uniformly over
$1\le L\le\ell_k$,
\[
 \frac{B_k(L)}{A_k}
 =\frac{x(x^L-1)}{k(x-1)}\{1+o(1)\}.
\]
The finite-radius factor $1-y_k^{-L}$ is deleted only when $L\to\infty$ as well as
$L=o(k)$.  In that regime the sharp transition variable is
$\delta_k=L-\log_xk$: the ratio tends to zero, a finite nonzero limit, or
infinity according as $\delta_k\to-\infty$, $c\in\mathbb R$, or
$+\infty$.  The finite limit is $x^{c+1}/(x-1)$.

Integer depth retains a real phase.  For
$L_k=\lfloor\log_xk+c\rfloor$ and
$\theta_k=\{\log_xk+c\}$,
\[
 \frac{B_k(L_k)}{A_k}
 =\frac{x^{c+1-\theta_k}}{x-1}\{1+o(1)\}.
\]
The phase limit set is $[0,1]$, so the ratio has no single limit; its
liminf and limsup are $x^c/(x-1)$ and $x^{c+1}/(x-1)$.  On the physical
noise clock the crossover lies only
$2L=(2/\log x)\log\log(1/\sigma)+O(1)$ orders above the first alias.
Actual-head statements remain conditional on the original unnormalized
same-clock hypothesis $D_{4k}(R)\to0$, which is not proved.  Linear depth,
direct/full-trace transfer, RH-288, Gates A--E, Hilbert--Polya, zero
identification, and RH remain open.
\end{abstract}

\section{Exact first-alias and post-alias ledgers}

Fix the physical constants
\begin{equation}\label{eq:constants}
 r_H=\frac{17}{20},\qquad R=\frac75,
 \qquad \frac{28}{17}<\lambda<\frac{17}{10}.
\end{equation}
On the physical first-alias clock $k=k_\sigma\to\infty$, RH-17 proves
\begin{equation}\label{eq:multiplier}
 |M_k|=C_M\lambda^k\{1+o(1)\},\qquad C_M>0.
\end{equation}
Put
\begin{equation}\label{eq:beta}
 \beta_k=\frac{|M_k|^{-1/(2k)}}{r_H},\qquad
 \beta=\frac1{r_H\sqrt\lambda},
\end{equation}
and
\begin{equation}\label{eq:x-y}
 x=(\beta R)^2=\frac{(R/r_H)^2}{\lambda}>1,
 \qquad
 y_k=(\beta_kR)^2
 =x\exp\left[-\frac{\log C_M}{k}+o(k^{-1})\right].
\end{equation}
These are the multiplier and Hardy-normalization locks used in RH-355
\citep{WangTimeOrdered2026,WangProjectorMass2026,WangUpperBurden2026}.

The finite graded counterloop is
\begin{equation}\label{eq:Y}
 \cY_k=\{\beta_ke^{ij\pi/k},\beta_ke^{-ij\pi/k}:1\le j\le k-1\}.
\end{equation}
Its exact power ledger is
\begin{equation}\label{eq:s-ledger}
 s_{k,n}=\sum_{\nu\in\cY_k}\nu^n
 =\beta_k^n\bigl(2k\one_{2k\mid n}-1-(-1)^n\bigr)
\end{equation}
\citep{WangAliasLedger2026,WangHeadCounterloop2026}.  Thus the first alias
at $n=2k$ has weighted modulus
\begin{equation}\label{eq:A}
 A_k:=\frac{|s_{k,2k}|R^{2k}}{2k}
 =\left(1-\frac1k\right)y_k^k.
\end{equation}
For an integer depth
\begin{equation}\label{eq:L-domain}
 1\le L\le k-1,
\end{equation}
define the even post-first-alias budget through $n=2k+2L$ by
\begin{equation}\label{eq:B}
 B_k(L):=\sum_{j=1}^L
 \frac{|s_{k,2k+2j}|R^{2k+2j}}{2k+2j}
 =\sum_{j=1}^L\frac{y_k^{k+j}}{k+j}.
\end{equation}
The upper limit in \eqref{eq:L-domain} keeps every order strictly below the
second alias $4k$.  Odd counterloop moments in this band vanish.

\begin{proposition}[Exact finite ratio]\label{prop:exact-ratio}
For every finite $k\ge2$ and every depth in \eqref{eq:L-domain},
\begin{equation}\label{eq:exact-ratio}
 \boxed{
 \frac{B_k(L)}{A_k}
 =\frac{k}{k-1}\sum_{j=1}^L\frac{y_k^j}{k+j}.}
\end{equation}
\end{proposition}

\begin{proof}
In the strict upper band, \eqref{eq:s-ledger} gives
$s_{k,2k+2j}=-2\beta_k^{2k+2j}$.  Hence the displayed summands in
\eqref{eq:B} are exact.  Divide their sum by \eqref{eq:A}.
\end{proof}

The common factor $y_k^k$ cancels exactly in
\eqref{eq:exact-ratio}.  The remaining finite-radius drift is carried by
$y_k^j$ and is negligible precisely when the post-alias depth is
sublinear.

\section{Uniform mesoscopic law}

\begin{theorem}[Sharp uniform post-alias law]\label{thm:uniform}
Let $\ell_k$ be any positive integer sequence with $\ell_k=o(k)$.  Then
\begin{equation}\label{eq:uniform-sup}
 \sup_{1\le L\le\ell_k}
 \left|
 \frac{B_k(L)/A_k}
 {x(x^L-1)/\{k(x-1)\}}-1
 \right|\longrightarrow0.
\end{equation}
Equivalently, uniformly on that domain,
\begin{equation}\label{eq:uniform-law}
 \boxed{
 \frac{B_k(L)}{A_k}
 =\frac{x(x^L-1)}{k(x-1)}\{1+o(1)\}.}
\end{equation}
The constant $C_M$ does not occur in the mesoscopic leading term.
\end{theorem}

\begin{proof}
Write
\begin{equation}\label{eq:epsilon}
 \varepsilon_k=\log(y_k/x)
 =-\frac{\log C_M}{k}+o(k^{-1}).
\end{equation}
Since $\ell_k=o(k)$,
\begin{equation}\label{eq:y-x-uniform}
 \sup_{1\le j\le\ell_k}
 \left|\frac{y_k^j}{x^j}-1\right|
 =\sup_{1\le j\le\ell_k}|e^{j\varepsilon_k}-1|
 \longrightarrow0.
\end{equation}
Also
\begin{equation}\label{eq:denominator-uniform}
 \sup_{1\le j\le\ell_k}
 \left|\frac{k^2}{(k-1)(k+j)}-1\right|
 \longrightarrow0.
\end{equation}
All summands are positive.  Therefore \eqref{eq:exact-ratio},
\eqref{eq:y-x-uniform}, and \eqref{eq:denominator-uniform} show, uniformly
in $L$, that
\begin{align}
 \frac{B_k(L)}{A_k}
 &=\frac{y_k(y_k^L-1)}{k(y_k-1)}\{1+o(1)\}
 \label{eq:y-geometric}\\
 &=\frac{y_k^{L+1}}{k(y_k-1)}
   (1-y_k^{-L})\{1+o(1)\}
 \label{eq:y-minus-one}\\
 &=\frac{1+o(1)}{k}\sum_{j=1}^Lx^j
 =\frac{x(x^L-1)}{k(x-1)}\{1+o(1)\}.
\end{align}
This proves \eqref{eq:uniform-sup}.  The cancellation of $C_M$ follows from
\eqref{eq:y-x-uniform}; no numerical value of $C_M$ is used.
\end{proof}

The geometric subtraction in \eqref{eq:uniform-law} must be retained at
bounded depth.

\begin{corollary}[Fixed and growing mesoscopic depths]
\label{cor:fixed-growing}
For every fixed integer $L\ge1$,
\begin{equation}\label{eq:fixed-L}
 k\frac{B_k(L)}{A_k}
 \longrightarrow\frac{x(x^L-1)}{x-1}.
\end{equation}
If instead
\begin{equation}\label{eq:growing-L}
 L\longrightarrow\infty,
 \qquad L=o(k),
\end{equation}
then and only at this quantified simplification step,
\begin{equation}\label{eq:master}
 \boxed{
 \frac{B_k(L)}{A_k}
 =\frac{x}{x-1}\,x^{L-\log_xk}\{1+o(1)\}.}
\end{equation}
\end{corollary}

\begin{proof}
Equation \eqref{eq:fixed-L} is the fixed-depth specialization of
\cref{thm:uniform}.  Under \eqref{eq:growing-L},
$y_k$ is eventually bounded below by a fixed number greater than one, so
$y_k^{-L}\to0$.  Thus and only under this additional $L\to\infty$
quantifier may the factor $1-y_k^{-L}$ in \eqref{eq:y-minus-one} be
deleted.  Equivalently,
$x^L-1=x^L(1-x^{-L})=x^L\{1+o(1)\}$; substitution in
\eqref{eq:uniform-law} gives \eqref{eq:master}.
\end{proof}

In particular, deleting $1-y_k^{-L}$ (or its limiting counterpart
$1-x^{-L}$) at fixed $L$ would change the leading
constant.  No finite table or informal use of ``mesoscopic'' licenses that
deletion.

\section{The crossover and its integer phase}

\begin{theorem}[Post-first-alias crossover]\label{thm:crossover}
Let $1\le L\le\ell_k$ for some $\ell_k=o(k)$, and put
\begin{equation}\label{eq:delta}
 \delta_k=L-\log_xk.
\end{equation}
Then:
\begin{enumerate}
 \item if $\delta_k\to-\infty$, then $B_k(L)/A_k\to0$;
 \item if $\delta_k\to c\in\mathbb R$, then
 \begin{equation}\label{eq:finite-crossover}
  \frac{B_k(L)}{A_k}\longrightarrow
  \frac{x^{c+1}}{x-1};
 \end{equation}
 \item if $\delta_k\to+\infty$, then $B_k(L)/A_k\to\infty$.
\end{enumerate}
The continuous balance offset is
\begin{equation}\label{eq:balance-offset}
 L=\log_xk+\log_x\left(\frac{x-1}{x}\right)+o(1).
\end{equation}
\end{theorem}

\begin{proof}
In the first case, \eqref{eq:uniform-law} and $x^L-1\le x^L$ give
\[
 0\le\frac{B_k(L)}{A_k}
 \le\frac{x}{x-1}x^{\delta_k}\{1+o(1)\}\longrightarrow0.
\]
In the finite and supercritical cases, $L\to\infty$ automatically, so
\eqref{eq:master} applies and gives the asserted limits.  Setting the finite
limit equal to one yields \eqref{eq:balance-offset}.
\end{proof}

Because the depth is integer, a fixed offset does not normally produce a
single limit.

\begin{theorem}[Floor-phase law and complete limit set]
\label{thm:integer-phase}
Fix $c\in\mathbb R$ and define
\begin{equation}\label{eq:floor-phase}
 L_k(c)=\lfloor\log_xk+c\rfloor,
 \qquad
 \theta_k(c)=\{\log_xk+c\}.
\end{equation}
Then
\begin{equation}\label{eq:phase-law}
 \boxed{
 \frac{B_k(L_k(c))}{A_k}
 =\frac{x^{c+1-\theta_k(c)}}{x-1}\{1+o(1)\}.}
\end{equation}
The limit set of $\theta_k(c)$ is the closed interval $[0,1]$.  Hence
\begin{align}
 \liminf_{k\to\infty}\frac{B_k(L_k(c))}{A_k}
 &=\frac{x^c}{x-1},
 \label{eq:liminf}\\
 \limsup_{k\to\infty}\frac{B_k(L_k(c))}{A_k}
 &=\frac{x^{c+1}}{x-1}.
 \label{eq:limsup}
\end{align}
In particular, the full sequence has no single limit.
\end{theorem}

\begin{proof}
Here $L_k(c)\sim\log_xk$, so $L_k(c)\to\infty$ and $L_k(c)=o(k)$.
Moreover
$L_k(c)-\log_xk=c-\theta_k(c)$.  Equation \eqref{eq:phase-law} follows
from \eqref{eq:master}.

For $t\in(0,1)$, take
$k_m=\lfloor x^{m+t-c}\rfloor$.  Then
$\log_xk_m+c=m+t+o(1)$, so a subsequence of phases tends to $t$.
For the endpoints, the choices
\begin{equation}\label{eq:phase-endpoint-subsequences}
 k_m^{(0)}=\left\lceil x^{m-c}\right\rceil,
 \qquad
 k_m^{(1)}=\left\lceil x^{m+1-c}\right\rceil-1
\end{equation}
give phases tending to zero and one, respectively.  The second formula also
handles the case in which $x^{m+1-c}$ is an integer: subtracting one keeps
the logarithm just below the next integer.  Thus the complete phase limit
set is $[0,1]$.  Apply the continuous decreasing function
$t\mapsto x^{c+1-t}/(x-1)$ to obtain \eqref{eq:liminf} and
\eqref{eq:limsup}.  Since $x>1$, the endpoints differ.
\end{proof}

On the physical clock
\begin{equation}\label{eq:physical-clock}
 k=\frac{\log(1/\sigma)}{2\log\lambda}+O(1).
\end{equation}
Therefore every bounded crossover window $L=\log_xk+O(1)$ lies at the
order displacement
\begin{equation}\label{eq:physical-depth}
 n-2k=2L
 =\frac{2}{\log x}\log\log(1/\sigma)+O(1).
\end{equation}
This is a post-first-alias mesoscopic scale: it diverges, but is
$o(k)$.

\section{Conditional actual-head transfer}

Let $\cH_\sigma$ be the actual modulus-complete Hardy head of RH-342, with
moments $h_{\sigma,n}$, and set
\begin{equation}\label{eq:defect}
 d_{\sigma,k,n}=h_{\sigma,n}-s_{k,n}.
\end{equation}
The original same-clock unnormalized transport leaf is
\begin{equation}\label{eq:D4k}
 D_{4k}(R)=\sum_{2\le n<4k}
 \frac{|d_{\sigma,k,n}|R^n}{n}\longrightarrow0
\end{equation}
\citep{WangGluing2026,WangSynchronization2026}.  It is an open hypothesis,
not a result of this paper.

Define the actual first-alias and even post-alias budgets
\begin{align}
 A_k^{\cH}
 &=\frac{|h_{\sigma,2k}|R^{2k}}{2k},
 \label{eq:AH}\\
 B_k^{\cH}(L)
 &=\sum_{j=1}^L
 \frac{|h_{\sigma,2k+2j}|R^{2k+2j}}{2k+2j},
 \label{eq:BH}
\end{align}
and let $O_k^{\cH}(L)$ be the corresponding weighted sum over odd orders
strictly between $2k$ and $2k+2L$.

\begin{theorem}[Conditional inheritance by the actual head]
\label{thm:actual-conditional}
Assume \eqref{eq:D4k} on one physical clock.  Then
\begin{equation}\label{eq:actual-alias}
 \frac{A_k^{\cH}}{A_k}\longrightarrow1,
 \qquad
 \sup_{1\le L\le k-1}
 \left|\frac{B_k^{\cH}(L)}{B_k(L)}-1\right|
 \longrightarrow0,
\end{equation}
and
\begin{equation}\label{eq:actual-odd}
 \sup_{1\le L\le k-1}O_k^{\cH}(L)\longrightarrow0.
\end{equation}
Consequently every mesoscopic conclusion of
\cref{thm:uniform,thm:crossover,thm:integer-phase} transfers to
$B_k^{\cH}(L)/A_k^{\cH}$ under this hypothesis.
\end{theorem}

\begin{proof}
The reverse triangle inequality gives
\begin{equation}\label{eq:actual-differences}
 |A_k^{\cH}-A_k|\le D_{4k}(R),
 \qquad
 |B_k^{\cH}(L)-B_k(L)|\le D_{4k}(R).
\end{equation}
Now $A_k\sim C_M^{-1}x^k\to\infty$, while uniformly in $L\ge1$,
\[
 B_k(L)\ge B_k(1)
 =\frac{y_k^{k+1}}{k+1}
 \sim\frac{x^{k+1}}{C_Mk}\longrightarrow\infty.
\]
Divide \eqref{eq:actual-differences} by these lower bounds and use
$D_{4k}(R)=o(1)$.  Odd counterloop moments vanish, so every
$O_k^{\cH}(L)$ is bounded by $D_{4k}(R)$.  Finally divide the two actual
budget asymptotics to transfer the ratios.
\end{proof}

This theorem is conditional only.  It identifies neither the actual head
rank nor its roots with the counterloop and does not prove
\eqref{eq:D4k}.

\section{Boundary and reproducibility}

The condition $L=o(k)$ is structural.  At a linear depth $L\sim\alpha k$,
the two uniform replacements used in \eqref{eq:y-x-uniform} and
\eqref{eq:denominator-uniform} fail: $y_k^L/x^L$ can retain a
$C_M^{-\alpha}$ factor, and $(k+j)^{-1}$ is not uniformly $k^{-1}$.
No linear-depth extension is claimed.  The exact ratio
\eqref{eq:exact-ratio} does, however, leave a source-backed next problem:
for $L/k\to\alpha\in(0,1]$, endpoint geometric domination should be audited
with the new factors $C_M^{-\alpha}$ and $(1+\alpha)^{-1}$.  This separate
profile is not proved here.

RH-354 controls the direct coefficient
\begin{equation}\label{eq:p-distinction}
 p_{\sigma,k,n}=\tau_{\sigma,n}-a_n
 =q_{\sigma,k,n}-d_{\sigma,k,n}.
\end{equation}
Neither the unconditional counterloop crossover nor the conditional head
statement proves a theorem for $p$, $q$, or the full-trace $E_{\rm off}$
budget \citep{WangDirectTail2026}.  The remaining typed gluing hypotheses
are still required.

The executable artifact checks the exact rational ledger at
$\lambda=5/3$, $C_M=1$, the full geometric factor $x^L-1$, and synthetic
finite-radius phase rows at archived high-precision diagnostics.  The
finite ratio itself is not required to lie inside its limiting phase
cluster at any prescribed $k$; only convergence to the phase law is tested.
The decimals are not interval certificates, actual noisy-head observations,
or asymptotic evidence.

RH-356 proves an unconditional mesoscopic crossover only for the graded
counterloop and a conditional transfer statement for the actual head.
It does not prove $D_{4k}(R)\to0$, root or rank matching, a linear-depth
law, direct/full-trace closure, the RH-241 moving noisy envelope, or RH-288
activation.  Gates A--E remain false/open.  No Hilbert--Polya operator,
Riemann-zero identification, von Mangoldt trace, completed-zeta divisor
equality, or proof of the Riemann Hypothesis is claimed.

\bibliographystyle{plainnat}
\bibliography{references}

\end{document}
